% !TeX root = ../induk.tex
\begin{table}[H]
\centering
\caption{Matriks Transisi Klaster Tahun 2021 ke 2022} \label{tab:transisi_2021_2022}
\begin{tabular}{|c|p{2cm}|p{2cm}|p{2cm}|}
\hline
\makecell{\textbf{Tahun 2021} \\ \textbf{2022}} & \textbf{Klaster 1 (2022)} & \textbf{Klaster 2 (2022)} & \textbf{Klaster 3 (2022)} \\ \hline
\textbf{Klaster 1 (2021)} & Sukabumi, Cianjur, Garut, Tasikmalaya, Ciamis, Kuningan, Cirebon, Majalengka, Indramayu, Subang, Bandung Barat, Pangandaran, Pandeglang, Lebak & - & Sumedang, Kota Banjar \\ \hline
\textbf{Klaster 2 (2021)} & - & Kota Jakarta Selatan, Kota Jakarta Timur, Kota Jakarta Pusat, Kota Jakarta Barat, Kota Jakarta Utara, Kota Bandung, Kota Bekasi, Kota Depok, Kota Cimahi, Kota Tangerang, Kota Tangerang Selatan & Kota Bogor \\ \hline
\textbf{Klaster 3 (2021)} & Purwakarta & - & Kep. Seribu, Bogor, Bandung, Karawang, Bekasi, Kota Sukabumi, Kota Cirebon, Kota Tasikmalaya, Tangerang, Serang, Kota Cilegon, Kota Serang \\ \hline
\end{tabular}
\end{table}